% =========================
% Cas pratique 6
% =========================

\section[Cas 6]{Cas pratique 6 : Droit conférés par une demande de brevet}

%--------------------------
\begin{frame}{Question 1}

\begin{block}{}
L'article \textbf{L613-1} CPI dispose que le droit exclusif d'exploitation mentionné 
à l'article L. 611-1 prend effet à compter du dépôt de la demande.\\
+ \textbf{L615-4} (opposable à partir de la publication)
+ \textbf{L615-3} (interdiction provisoire) "toute personne habilitée
à agir en contrefaçon".
\end{block}

\begin{itemize}
  \item En l'espèce, la demande de brevet EP n'a pas été délivrée, mais
  la demande de brevet EP a été publiée.
  \item En conclusion, la société Hélicopathe est recevable pour agir
  en contrefaçon. La société hélicopathe a le droit de demander
  une interdiction provisoire.
\end{itemize}

\begin{block}{}
L'article \textbf{62} AJUB dispose que la Juridiction peut
prononcer des injections visant à 
à interdire, à titre provisoire et
sous réserve, le cas échéant, du paie
ment d'une astreinte, que la contrefaçon
présumée se poursuive.
\end{block}

\end{frame}

%--------------------------
\begin{frame}{Question 2}

\begin{itemize}
  \item En l'espèce, la demande de brevet EP
  la portée de la demande de brevet EP a été réduite.
  \item En conclusion, les chances délivrances d'un
  brevet EP sont plus grandes mais risque que le produit
  contrefaisant "sorte" de la portée du brevet EP délivré.
\end{itemize}

\end{frame}

%--------------------------
\begin{frame}{Question 3}

\begin{block}{}
L'article \textbf{L615-4} dispose que
par exception aux dispositions de l'article L. 613-1,
 les faits antérieurs à la date à laquelle
  la demande de brevet a été rendue publique 
  en vertu de l'article L. 612-21 ou
   à celle de la notification à tout tiers 
   d'une copie certifiée de cette demande 
   ne sont pas considérés comme 
ayant porté atteinte aux droits attachés au brevet. 
\end{block}

\begin{itemize}
  \item En l'espèce, la demande de brevet EP
  a été publiée.
  \item En conclusion, les actes de contrefaçon sont
  comptabilisés à partir de la date de publication
  de la demande de brevet EP.
\end{itemize}

\end{frame}

%--------------------------
\begin{frame}{Question 4}

\begin{block}{}
\begin{itemize}
    \item TGI de Paris: Article \textbf{378 CPC} (sursis pour bonne administration de justice).
    \item JUB: Article \textbf{33}(10) AJUB dispose que
les parties informent la Juridiction
de toute procédure de nullité, de limitation 
ou d'opposition pendante devant
l'Office européen des brevets, ainsi que
de toute demande de procédure accélérée
 présentée auprès de l'Office européen des brevets.
  La Juridiction peut
suspendre la procédure lorsqu'une
décision rapide peut être attendue de
l'Office européen des brevets.
\end{itemize}
\end{block}

\begin{itemize}
    \item En conclusion, devant TGI de Paris,
    un sursis à statuer est à prévoir.
    Devant la JUB, le sursis à statuer
    dépend de l'avancement de l'opposition.
\end{itemize}

\end{frame}

